% Curriculum Vitae -- M. Harper Langston
%
% Build:   python3 build_lists.py && latexmk -pdf cv.tex
%
% Publications, preprints/reports and patents are generated from the website
% (../index.html) by build_lists.py, so add new papers there and re-run it.
% Everything else is edited directly in this file.
%
% Items marked \todo{...} print in red: fill them in, then delete the \todo.

\documentclass[10pt,letterpaper]{article}

\usepackage[T1]{fontenc}
\usepackage[utf8]{inputenc}
\usepackage{mathpazo}                 % Palatino, to match the website's headings
\usepackage[scaled=0.9]{helvet}
\usepackage{microtype}
\usepackage[margin=0.85in,top=0.75in,bottom=0.8in]{geometry}
\usepackage{xcolor}
\usepackage{enumitem}
\usepackage{titlesec}
\usepackage{fancyhdr}
\usepackage{lastpage}
\usepackage[hidelinks]{hyperref}

\definecolor{accent}{HTML}{1F5F8B}
\definecolor{muted}{HTML}{5C6370}
\definecolor{award}{HTML}{8A5A00}

\hypersetup{pdftitle={M. Harper Langston -- Curriculum Vitae}, pdfauthor={M. Harper Langston}}

% ---------- layout helpers ----------
\setlength{\parindent}{0pt}
\setlength{\parskip}{0pt}
\pagestyle{fancy}
\fancyhf{}
\renewcommand{\headrulewidth}{0pt}
\fancyfoot[L]{\footnotesize\color{muted}M. Harper Langston \textperiodcentered{} Curriculum Vitae}
\fancyfoot[R]{\footnotesize\color{muted}\thepage\ of \pageref{LastPage}}

\titleformat{\section}{\large\bfseries\color{accent}}{}{0pt}{}[\vspace{1pt}{\color{accent!40}\titlerule[0.6pt]}]
\titlespacing*{\section}{0pt}{14pt}{6pt}

\newcommand{\todo}[1]{{\color{red}\textbf{[#1]}}}

% \job{Organization}{Location}{Title}{Dates}{description}
\newcommand{\job}[5]{%
  \textbf{#1}\hfill{\color{muted}#2}\\
  \textit{#3}\hfill{\color{muted}#4}%
  \ifx&#5&\else\\[1pt]{\small #5}\fi
  \par\vspace{6pt}}

% \role{Title}{Dates} -- an extra role at the same organization
\newcommand{\role}[2]{\textit{#1}\hfill{\color{muted}#2}\par\vspace{2pt}}

% Publication lists
\newcommand{\yearhead}[1]{\par\vspace{5pt}{\small\bfseries\color{muted}#1}\par\vspace{2pt}}
\newcommand{\subhead}[1]{\par\vspace{6pt}{\bfseries #1}\par\vspace{3pt}}
\newcommand{\award}[1]{\,{\color{award}\textbf{(#1)}}}
\newcounter{pubno}
\newcommand{\pub}[4]{%
  \stepcounter{pubno}%
  \par\noindent\hangindent=2.4em\hangafter=1
  \makebox[2.4em][l]{\small\color{muted}[\thepubno]}%
  \textbf{#1}. {\small #2}. \textit{\small #3}.#4\par\vspace{3pt}}
\newcommand{\resetpubs}[1]{\setcounter{pubno}{0}\renewcommand{\thepubno}{#1\arabic{pubno}}}

\setlist[itemize]{leftmargin=1.2em,itemsep=1pt,topsep=2pt}

% ---------- document ----------
\begin{document}

{\centering
  {\fontsize{24}{28}\selectfont\bfseries M. Harper Langston}\\[4pt]
  {\large Principal Engineer, Qualcomm AI Research}\\[6pt]
  {\small\color{muted}
    \href{https://harperlangston.com}{harperlangston.com} \enspace\textperiodcentered\enspace
    \href{mailto:harperlangston@gmail.com}{harperlangston@gmail.com} \enspace\textperiodcentered\enspace
    \href{https://github.com/harperlangston}{github.com/harperlangston} \enspace\textperiodcentered\enspace
    \href{https://scholar.google.com/citations?user=I1-ydOQAAAAJ}{Google Scholar}}\par}

\section{Research Interests}
Efficient machine learning (attention and KV-cache compression, polynomial and linear attention, hardware-aware
inference), global polynomial and combinatorial optimization, fast multipole methods and fast solvers,
preconditioners and multigrid, tensor decompositions, and high-performance and parallel computing.

\section{Experience}

\job{Qualcomm AI Research, Qualcomm Technologies, Inc.}{\todo{location}}
    {Principal Engineer}{Nov 2021 -- Present}
    {Lead research and engineering efforts spanning efficient machine learning systems, large language model
     inference, optimization research and design, fast algorithms, hardware-aware AI and next-generation neural
     architectures, focused on co-designing algorithms, systems and hardware for scalable, efficient AI. Recent work
     includes KV-cache selection and eviction for LLM inference (QUOKA, CAOTE), polynomial attention for vision
     transformers (PADRe), NPU-friendly linear attention, global polynomial and MAX-CUT optimization, and learned
     optimizer schedules.}

\job{Reservoir Labs, Inc.}{New York, NY}
    {Managing Engineer}{Feb 2015 -- Nov 2021}
    {Managed a team of engineers and researchers developing algorithms and software for signal processing,
     machine learning, fast solvers, fast Fourier transforms, tensor decompositions, fast polynomial solvers and
     radar. PI for DOE SBIR projects including MACH-B (Multipole Accelerator Codes for Hadron Beams) and PUMA-V,
     managing subcontractors at Stony Brook University and Brookhaven National Laboratory; managed additional
     SBIR and government projects and wrote funding proposals. Introduced code review and project tracking.
     Production and research software in C/C++, Java and Python.}
\vspace{-4pt}
\role{Senior Research Engineer}{Nov 2012 -- Jan 2015}
{\small HPC research for the exascale software stack; power-efficient tools for GPUs and distributed-memory CPUs;
 sparse tensor decompositions; graph sparsification and solvers for diagonally dominant systems; fast
 multidimensional solvers; the R-Stream auto-parallelizing compiler.}\par\vspace{6pt}

\job{Courant Institute of Mathematical Sciences, New York University}{New York, NY}
    {Adjunct Assistant Professor}{2013 -- \todo{end date or Present}}
    {Designed and taught graduate courses in the Mathematics and Computer Science departments.}

\job{Perceptive Pixel}{New York, NY and Washington, DC}
    {Software Engineer}{Jan 2012 -- Nov 2012}
    {Tools and software for multi-touch hardware: recording, playback and visualization of touch-sensor input,
     synchronized video capture, and multi-user document manipulation; supported CNN applications during the
     2012 election season. C++.}

\job{Courant Institute of Mathematical Sciences, New York University}{New York, NY}
    {Research Assistant}{Sep 2003 -- Aug 2008; Sep 2010 -- Jan 2012}
    {Fast, parallel, scalable 3D PDE solvers in complex geometries with non-homogeneous forces; adaptive particle
     and vortex methods; visualization of embedded-boundary fluid simulations. Supported by DOE and PPPL.}

\job{Center for Plasma Edge Simulation}{New York, NY and Washington, DC}
    {Research Scientist}{Jul 2010 -- Sep 2010}
    {FMM-based cross-verification tools for plasma edge simulation codes in 3D toroidal geometries (DOE, PPPL).}

\job{Georgia Institute of Technology}{Atlanta, GA}
    {Computational Science and Engineering Research Scientist}{Sep 2008 -- Apr 2010}
    {Brain-tumor growth modeling for predictive medical imaging: multi-threaded diffusion and multigrid elasticity
     solvers with parallel octree registration (NIH, DoD; with the University of Pennsylvania). Research Affiliate,
     May 2010 -- May 2011.}

\textbf{Earlier positions}\par\vspace{2pt}
{\small
\begin{itemize}
  \item Junior Research Scientist, NYU Media Research Lab, Nov 2000 -- Aug 2002 (teleimmersive responsive workbench).
  \item Webmaster, NYU Computer Science Department, Dec 2000 -- May 2002.
  \item Software Engineering Contractor, Jaron Lanier / Advanced Network \& Services, Sep -- Nov 2001.
  \item Operations/Database Developer, Mighty Seven Networks, Jan -- Sep 2001.
  \item Innovation and Technology Management Research Specialist, McKinsey \& Company, Oct 1997 -- Nov 2000.
  \item Mathematics Department Research Assistant, Bowdoin College, Jan -- May 1997 (with Wells Johnson).
\end{itemize}}

\section{Education}

\job{New York University, Courant Institute of Mathematical Sciences}{New York, NY}
    {Ph.D., Computer Science}{January 2012}
    {Dissertation: \textit{An Adaptive Fast Volume Solver in Three Dimensions for Free-Space, Periodic and
     Dirichlet Boundary Conditions}. Advisors: Denis Zorin and Leslie Greengard. McCracken Fellowship.}

\job{New York University, Courant Institute of Mathematical Sciences}{New York, NY}
    {M.S., Scientific Computing}{May 2003}
    {Thesis: \textit{Investigation of Computational and Visualization Methods for the Incompressible
     Navier--Stokes Equations}. Advisors: Denis Zorin and George Biros; academic advisor: Michael Shelley.}

\job{Bowdoin College}{Brunswick, ME}
    {B.A., Mathematics and Russian Language, Magna Cum Laude}{June 1997}
    {Smyth Mathematics Prize (1995); James Bowdoin Scholar (1994--1997). Also: Middlebury College Intensive
     Russian Summer Program (1994).}

\section{Honors and Awards}
\begin{itemize}
  \item Best Paper Award, IEEE High Performance Extreme Computing Conference (HPEC), 2017.
  \item Best Paper Finalist, ACM/IEEE Supercomputing (SC09), 2009.
  \item Best Student Paper and Gordon Bell Prize Finalist, ACM/IEEE Supercomputing (SC03), 2003.
  \item McCracken Fellowship, New York University.
  \item Smyth Mathematics Prize, Bowdoin College, 1995; James Bowdoin Scholar, 1994--1997.
  \item \todo{add recent honors}
\end{itemize}

\section{Publications}
{\small\color{muted}Conference, journal and workshop papers, newest first.}\par
% Generated by build_lists.py from ../index.html -- edit the website, not this file.

\yearhead{2026}
\pub{\href{https://arxiv.org/abs/2610.19043}{Graph Coarsening Enables Escaping Local Optima in Compute-Constrained Nonconvex MAX-CUT}}{\textbf{M. Harper Langston}, Pierre-David Letourneau, Dalton Jones, Harris Teague, Peter Scott, Mingu Lee, Richard A. Lethin, Christopher Lott, Roland Memisevic, Fatih Porikli}{OPT 2026: Optimization for Machine Learning, NeurIPS 2026 Workshop (Sydney)}{}
\pub{\href{https://arxiv.org/abs/2610.32801}{Multiplicative vs. Additive State Coupling in Learned Optimizer Schedules}}{\textbf{M. Harper Langston}, Pierre-David Letourneau, Dalton Jones, Roland Memisevic, Peter Scott, Mingu Lee, Richard A. Lethin, Christopher Lott, Fatih Porikli, Harris Teague}{OPT 2026: Optimization for Machine Learning, NeurIPS 2026 Workshop (Sydney)}{}
\pub{\href{https://arxiv.org/abs/2602.08722}{QUOKA: Query-Oriented KV Selection for Efficient LLM Prefill}}{Dalton Jones, Junyoung Park, Matthew Morse, Mingu Lee, Chris Lott, \textbf{M. Harper Langston}}{International Conference on Learning Representations (ICLR), 2026}{}
\pub{\href{https://arxiv.org/abs/2504.14051}{CAOTE: Optimizing KV Cache Memory Through Attention Output Error-based Token Eviction}}{Raghavv Goel, Junyoung Park, Mukul Gagrani, Dalton Jones, Matthew Morse, \textbf{M. Harper Langston}, Christopher Lott, Mingu Lee}{ICLR 2026 Workshop on Memory for LLM-Based Agentic Systems (MemAgents), 2026}{}
\pub{\href{https://arxiv.org/abs/2606.06034}{When Good Enough Is Optimal: Multiplication-Only Matrix Inversion Approximation for Quantized Gated DeltaNet}}{Luoming Zhang, Yuwei Ren, Kui Zhang, Tian Liu, Lingjuan Ge, Denghao Li, \textbf{M. Harper Langston}, Yin Huang, Weiliang Will Zeng, Liang Zhang}{International Conference on Machine Learning (ICML), 2026}{}
\yearhead{2025}
\pub{\href{https://arxiv.org/abs/2407.11306}{PADRe: A Unifying Polynomial Attention Drop-in Replacement for Efficient Vision Transformer}}{Pierre-David Letourneau, Manish Kumar Singh, Hsin-Pai Cheng, Shizhong Han, Yunxiao Shi, Dalton Jones, \textbf{M. Harper Langston}, Hong Cai, Fatih Porikli}{International Conference on Learning Representations (ICLR), 2025}{}
\yearhead{2022}
\pub{\href{https://doi.org/10.18429/JACoW-ICALEPCS2021-MOPV048}{Fast Multipole Method (FMM)-Based Particle Accelerator Simulations in the Context of Tune Depression Studies}}{\textbf{M. Harper Langston}, Richard Lethin, Pierre-David Letourneau, Julia Wei}{18th International Conference on Accelerator and Large Experimental Physics Control Systems (ICALEPCS 2021), Shanghai, China, pp. 271--276, 2022}{}
\yearhead{2021}
\pub{\href{https://harperlangston.com/papers/Sagan\_2021\_J.\_Inst.\_16\_T10002.pdf}{Simulations of Future Particle Accelerators: Issues and Mitigations}}{D. Sagan, M. Berz, N.M. Cook, Y. Hao, G. Hoffstaetter, A. Huebl, C.-K. Huang, \textbf{M.H. Langston}, C.E. Mayes, C.E. Mitchell, C.-K. Ng, J. Qiang, R.D. Ryne, A. Scheinker, E. Stern, J.-L. Vay, D. Winklehner and H. Zhang}{Journal of Instrumentation (JINST), IOP Publishing, Volume 16, October 2021}{}
\pub{\href{https://harperlangston.com/papers/THPAB227.pdf}{MACH-B: Fast Multipole Method Approaches in Particle Accelerator Simulations for the Computational and Intensity Frontiers}}{\textbf{M. Harper Langston}, Julia Wei, Pierre-David Letourneau, Matthew J. Morse, Richard Lethin}{12th International Particle Accelerator Conference (IPAC'21), Campinas, SP, Brazil, 2021}{}
\pub{Boundary Integral Solver Approaches for Particle Accelerator Simulation Problems and Deployment on NERSC Hardware}{Julia Wei, \textbf{M. Harper Langston}, Pierre-David Letourneau, Matthew J. Morse, Larry Weintraub, Aimee Nogoy, Noah Amsel, Richard Lethin}{IEEE High Performance Extreme Computing Conference (HPEC), 2021}{}
\yearhead{2020}
\pub{\href{https://www.osti.gov/servlets/purl/1673301}{Approximate Inverse Chain Preconditioner: Iteration Count Case Study for Spectral Support Solvers}}{\textbf{M. Harper Langston}, Pierre-David Letourneau, Julia Wei, Larry Weintraub, Mitchell Harris, Richard Lethin, Eric Papenhausen, Meifeng Lin}{IEEE High Performance Extreme Computing Conference (HPEC) 2020}{}
\pub{\href{https://harperlangston.com/papers/HPEC-Multiscale-Data-Analysis-Using-Binning-Tensor-Decompositions-and-Backtracking.pdf}{Multiscale Data Analysis Using Binning, Tensor Decompositions, and Backtracking}}{Dimitri Leggas, Thomas S. Henretty, James Ezick, Muthu Baskaran, Brendan von Hofe, Grace Cimaszewski, \textbf{M. Harper Langston}, Richard Lethin}{IEEE High Performance Extreme Computing Conference (HPEC) 2020}{}
\pub{\href{https://doi.org/10.1137/19M1279502}{Low-Frequency Electromagnetic Imaging Using Sensitivity Functions and Beamforming}}{Pierre-David Letourneau, Mitchell Tong Harris, Matthew \textbf{Harper Langston}, and George Papanicolaou}{SIAM Journal on Imaging Sciences, Vol. 13, No. 2, pp. 807-843, 2020}{}
\pub{\href{https://harperlangston.com/papers/3366707.pdf}{On the Bottleneck Structure of Congestion-Controlled Networks}}{Jordi Ros-Giralt, Sruthi Yellamraju, Atul Bohara, \textbf{M. Harper Langston}, Richard Lethin, Yuang Jiang, Leandros Tassiulas, Josie Li, Yuang Tan, Malathi Veeraraghavan}{ACM SIGMETRICS: International Conference on Measurement and Modeling of Computer Systems, 2020}{}
\yearhead{2019}
\pub{\href{https://harperlangston.com/papers/Combinatorial\_Multigrid\_Advanced\_Preconditioners\_For\_Ill\_Conditioned\_Linear\_Systems.pdf}{Combinatorial Multigrid: Advanced Preconditioners For Ill-Conditioned Linear Systems}}{\textbf{M. Harper Langston}, Mitchell Tong Harris, Pierre-David Letourneau, Richard Lethin, James Ezick}{IEEE Conference on High Performance Extreme Computing (HPEC '19), 2019}{}
\pub{\href{https://harperlangston.com/papers/Fast-Large-Scale-Algorithm-for-Electromagnetic-Wave-Propagation-in-3D-Media.pdf}{Fast Large-Scale Algorithm for Electromagnetic Wave Propagation in 3D Media}}{Mitchell Tong Harris, \textbf{M. Harper Langston}, Pierre-David Letourneau, George Papanicolaou, James Ezick, Richard Lethin}{IEEE Conference on High Performance Extreme Computing (HPEC '19), 2019}{}
\pub{\href{https://harperlangston.com/papers/PUMA-V.pdf}{PUMA-V: Optimizing Parallel Code Performance Through Interactive Visualization}}{Eric Papenhausen, \textbf{M. Harper Langston}, Benoit Meister, Richard Lethin, Klaus Mueller}{IEEE Computer Graphics and Applications 39(1): 84-99, 2019}{}
\pub{\href{https://harperlangston.com/papers/G2\_INDIS.pdf}{G2: A Network Optimization Framework for High-Precision Analysis of Bottleneck and Flow Performance}}{Jordi Ros-Giralt, Sruthi Yellamraju, Atul Bohara, \textbf{Harper Langston}, Richard Lethin, Yuang Jiang, Leandros Tassiulas, Josie Li, Ying Lin, Yuanlong Tan, Malathi Veeraraghavan}{2019 IEEE/ACM Supercomputing Conference, Innovating the Network for Data-Intensive Science (INDIS) Workshop}{}
\pub{\href{https://harperlangston.com/papers/RL\_SIAM\_NS19\_Abstract.pdf}{On the Bottleneck Structure of Positive Linear Programming}}{Jordi Ros-Giralt, \textbf{Harper Langston}, Aditya Gudibanda, Richard Lethin}{2019 SIAM Workshop on Network Science}{}
\yearhead{2018}
\pub{\href{https://doi.org/10.1117/12.2306933}{Topic Modeling for Analysis of Big Data Tensor Decompositions}}{Thomas Henretty, \textbf{M. Harper Langston}, Muthu Baskaran, James Ezick, Richard Lethin}{SPIE Proceedings Vol. 10652, Disruptive Technologies in Information Sciences, 1065208, 2018}{}
\yearhead{2017}
\pub{\href{https://harperlangston.com/papers/Memory-Efficient-Parallel-Sparse-Tensor-Decomposition-HPEC-2017.pdf}{Memory-efficient Parallel Tensor Decompositions}}{Muthu Baskaran, Tom Henretty, Benoit Pradelle, \textbf{M. Harper Langston}, David Bruns-Smith, James Ezick, Richard Lethin}{IEEE High Performance Extreme Computing Conference (HPEC '17), Waltham, MA, 2017}{\award{Best Paper Award}}
\pub{Discovering Deep Patterns In Large Scale Network Flows Using Tensor Decompositions}{James Ezick, Muthu Baskaran, David Bruns-Smith, Alan Commike, Thomas Henretty, \textbf{M. Harper Langston}, Jordi Ros-Giralt, Richard Lethin}{FloCon 2017, San Diego, CA}{}
\yearhead{2016}
\pub{\href{https://harperlangston.com/papers/lowrank\_hpec.pdf}{Accelerated Low-rank Updates to Tensor Decompositions}}{Muthu Baskaran, \textbf{M. Harper Langston}, Tahina Ramananandro, David Bruns-Smith, Tom Henretty, James Ezick, Richard Lethin}{IEEE Conference on High Performance Extreme Computing (HPEC '16), 2016}{}
\pub{\href{https://harperlangston.com/papers/sparse\_fft\_hpec.pdf}{A Sparse Multi-Dimensional Fast Fourier Transform with Stability to Noise in the Context of Image Processing and Change Detection}}{Pierre-David Letourneau, \textbf{M. Harper Langston}, Richard Lethin}{IEEE Conference on High Performance Extreme Computing (HPEC '16), 2016}{}
\pub{\href{https://harperlangston.com/papers/PID4293569.pdf}{An Interactive Visual Tool for Code Optimization and Parallelization Based on the Polyhedral Model}}{Eric Papenhausen, Klaus Mueller, \textbf{M. Harper Langston}, Benoit Meister, Richard Lethin}{Sixth International Workshop on Parallel Software Tools and Tool Infrastructures (PSTI 2016), held with ICPP 2016, the 45th International Conference on Parallel Processing, Philadelphia, PA, August 16--19, 2016}{}
\pub{Optimization of the Domain Dslash Operator for Intel Xeon CPUs}{Meifeng Lin, Eric Papenhausen, \textbf{M. Harper Langston}, Benoit Meister, Muthu Baskaran, Chulwoo Jung, Taku Izubuchi}{34th International Symposium on Lattice Field Theory, 2016}{}
\pub{PUMA-V: An Interactive Visual Tool for Code Optimization and Parallelization Based on the Polyhedral Model}{Eric Papenhausen, Klaus Mueller, \textbf{M. Harper Langston}, Benoit Meister, Richard Lethin}{2016 New York Scientific Data Summit (NYSDS), pp. 1--4, 2016}{}
\yearhead{2015}
\pub{\href{https://harperlangston.com/papers/vissoft15.pdf}{Polyhedral User Mapping and Assistant Visualizer Tool for the R-Stream Auto-Parallelizing Compiler}}{Eric Papenhausen, Bing Wang, \textbf{Harper Langston}, Muthu Baskaran, Tom Henretty, Taku Izubuchi, Ann Johnson, Chulwoo Jung, Meifeng Lin, Benoit Meister, Klaus Mueller, Richard Lethin}{IEEE 3rd annual Working Conference on Software Visualization (IEEE VISSOFT), Bremen, Germany, 2015}{}
\pub{\href{https://arxiv.org/abs/1512.01542}{Optimizing the domain wall fermion Dirac operator using the R-Stream source-to-source compiler}}{Meifeng Lin, Eric Papenhausen, \textbf{M. Harper Langston}, Benoit Meister, Muthu Baskaran, Taku Izubuchi, Chulwoo Jung}{33rd International Symposium on Lattice Field Theory, Kobe International Conference Center, Kobe, Japan, July 14--18, 2015}{}
\yearhead{2013}
\pub{\href{https://harperlangston.com/papers/hpec13.pdf}{Re-Introduction of Communication-Avoiding FMM-Accelerated FFTs with GPU Acceleration}}{\textbf{M. Harper Langston}, Muthu Baskaran, Benoît Meister, Nicolas Vasilache and Richard Lethin}{IEEE Conference on High Performance Extreme Computing (HPEC '13), 2013}{}
\yearhead{2012}
\pub{\href{https://harperlangston.com/papers/ACM2012.pdf}{A massively parallel adaptive fast-multipole method on heterogeneous architectures}}{I. Lashuk, A. Chandramowlishwaran, \textbf{H. Langston}, T.-A. Nguyen, R. S. Sampath, A. Shringarpure, R. W. Vuduc, L. Ying, D. Zorin, and G. Biros}{Communications of the ACM, vol. 55, 5, 2012, pp. 101-109}{}
\yearhead{2011}
\pub{\href{https://harperlangston.com/papers/camcos2011.pdf}{A Free-Space Adaptive FMM-Based PDE Solver in Three Dimensions}}{\textbf{H. Langston}, L. Greengard, and D. Zorin}{Communications in Applied Mathematics and Computational Science, vol. 6, no. 1, 2011, pp. 79--122}{}
\yearhead{2009}
\pub{\href{https://harperlangston.com/papers/sc09.pdf}{A massively parallel adaptive fast-multipole method on heterogeneous architectures}}{I. Lashuk, A. Chandramowlishwaran, \textbf{H. Langston}, T.-A. Nguyen, R. S. Sampath, A. Shringarpure, R. W. Vuduc, L. Ying, D. Zorin, and G. Biros}{Proceedings of SC09, ACM/IEEE Conference on High Performance Computing, pp. 1--11, 2009}{\award{Best Paper Finalist}}
\yearhead{2003}
\pub{\href{https://harperlangston.com/papers/ying.pdf}{A New Parallel Kernel-Independent Fast Multipole Method}}{L. Ying, G. Biros, \textbf{H. Langston}, and D. Zorin}{Proceedings of SC03, ACM/IEEE Conference on High Performance Computing, pp. 1--11, 2003}{\award{Best Student Paper \textperiodcentered{} Gordon Bell Prize Finalist}}


\section{Preprints and Technical Reports}
\resetpubs{R}
% Generated by build_lists.py from ../index.html -- edit the website, not this file.

\pub{\href{https://arxiv.org/abs/2311.02037}{An Efficient Framework for Global Non-Convex Polynomial Optimization with Algebraic Constraints}}{Mitchell Tong Harris, Pierre-David Letourneau, Dalton Jones, \textbf{M. Harper Langston}}{Preprint, arXiv:2311.02037, 2023}{}
\pub{\href{https://arxiv.org/abs/2309.14584}{A Sparse Fast Chebyshev Transform for High-Dimensional Approximation}}{Dalton Jones, Pierre-David Letourneau, Matthew J. Morse, \textbf{M. Harper Langston}}{Preprint, arXiv:2309.14584, 2023}{}
\pub{\href{https://arxiv.org/abs/2308.16731}{An Efficient Framework for Global Non-Convex Polynomial Optimization over the Hypercube}}{Pierre-David Letourneau, Dalton Jones, Matthew Morse, \textbf{M. Harper Langston}}{Preprint, arXiv:2308.16731, 2023}{}
\pub{MACH-B: Multipole Accelerator Codes for Hadron Beams --- DOE Phase I Final Technical Report}{\textbf{M. Harper Langston}, Pierre-David Letourneau, Julia Wei, Richard Lethin, A. Ryan}{Technical Report, Reservoir Labs, 2021}{}
\pub{Correlating and Analyzing Network Data through Interpretable Decompositions (CANDID)}{M. Baskaran, J. Ezick, T. Henretty, B. von Hofe, D. Leggas, \textbf{M. H. Langston}, et al.}{Technical Report, Reservoir Labs, 2021}{}
\pub{Fast Multipole Method Approaches in Particle Accelerator Simulations for the Computational and Intensity Frontiers}{\textbf{M. Harper Langston}, Pierre-David Letourneau, Julia Wei, Richard Lethin}{Snowmass 2021 Letter of Interest}{}
\pub{Ensign}{M. Baskaran, T. Henretty, D. Leggas, B. von Hofe, \textbf{M. H. Langston}, J. Ezick, et al.}{Technical Report, Reservoir Labs, 2020}{}
\pub{Adaptive Inhomogeneous PDE Solvers for Complex Geometries}{\textbf{H. Langston}, L. Greengard, and D. Zorin}{In preparation}{}
\pub{\href{https://arxiv.org/abs/1604.06682}{A sparse multidimensional FFT for real positive vectors}}{Pierre-David Letourneau, \textbf{M. Harper Langston}, Benoit Meister, Richard Lethin}{Preprint, arXiv:1604.06682, 2016}{}
\pub{PUMA-V: Polyhedral User Mapping Assistant and Visualizer --- Final Technical Report}{\textbf{M. H. Langston}, M. Lin, E. Papenhausen, B. Meister, A. Konstantinidis, et al.}{Technical Report, Reservoir Labs, New York, NY, 2016}{}
\pub{\href{https://arxiv.org/abs/1409.1914}{A Tale of Three Runtimes}}{Nicolas Vasilache, Muthu Baskaran, Tom Henretty, Benoit Meister, \textbf{M. Harper Langston}, Sanket Tavarageri, Richard Lethin}{Preprint, arXiv:1409.1914, 2014}{}
\pub{\href{https://harperlangston.com/papers/ascac\_top\_ten\_exascale\_2014.pdf}{DOE Advanced Scientific Computing Advisory Subcommittee (ASCAC) Report: Top Ten Exascale Research Challenges}}{R. Lucas (chair), J. Ang, K. Bergman, S. Borkar, W. Carlson, L. Carrington, G. Chiu, et al. \textperiodcentered{} \textbf{H. Langston} (contributor)}{ASCAC Subcommittee Report, U.S. Department of Energy, Office of Science, February 10, 2014}{}
\pub{\href{https://harperlangston.com/papers/vortex\_tr.pdf}{Avoiding Particle Dissipation for Adaptive Vortex Methods Through Circulation Conservation}}{\textbf{H. Langston}}{Technical Report, 2006}{}
\pub{\href{https://harperlangston.com/papers/cash.pdf}{Cash: Distributed Cooperative Buffer Caching}}{C. Decoro, \textbf{H. Langston}, J. Weinberger}{Technical Report, 2004}{}


\section{Patents}
\resetpubs{P}
% Generated by build_lists.py from ../index.html -- edit the website, not this file.

\subhead{Granted}
\pub{\href{https://patents.google.com/patent/US12229216B2/}{Kinematic Invariant-Space Maximum Entropy Tracker (KISMET)}}{Lawrence Craig Weintraub, \textbf{Matthew Harper Langston}, Julia Wei, Richard Lethin, Aimee Kristine Nogoy, Mitchell Harris, Paul Mountcastle}{US 12,229,216 B2 \textperiodcentered{} Granted February 18, 2025}{}
\pub{\href{https://patents.google.com/patent/US11907549B2/}{Systems and Methods for Minimizing Communications}}{Muthu M. Baskaran, Thomas Henretty, Ann Johnson, Athanasios Konstantinidis, \textbf{M. H. Langston}, Janice O. McMahon, Benoit J. Meister, Paul D. Mountcastle, Aale Naqvi, Benoit Pradelle, Tahina Ramananandro, Sanket Tavarageri, Richard A. Lethin}{US 11,907,549 B2 \textperiodcentered{} Granted February 20, 2024}{}
\pub{\href{https://patents.google.com/patent/US11789769B2/}{System and Method for Generation of Event Driven, Tuple-Space Based Programs}}{Muthu M. Baskaran, Thomas Henretty, \textbf{M. H. Langston}, Richard A. Lethin, Benoit J. Meister, Nicolas T. Vasilache, David E. Wohlford}{US 11,789,769 B2 \textperiodcentered{} Granted October 17, 2023}{}
\pub{\href{https://patents.google.com/patent/US11726197B2/}{Systems and Methods for Efficient Targeting}}{Muthu M. Baskaran, Thomas Henretty, Ann Johnson, Athanasios Konstantinidis, \textbf{M. H. Langston}, Janice O. McMahon, Benoit J. Meister, Paul D. Mountcastle, Aale Naqvi, Benoit Pradelle, Tahina Ramananandro, Sanket Tavarageri, Richard A. Lethin}{US 11,726,197 B2 \textperiodcentered{} Granted August 15, 2023}{}
\pub{\href{https://patents.google.com/patent/US11500557B2/}{Systems and Methods for Energy Proportional Scheduling}}{Muthu M. Baskaran, Thomas Henretty, Ann Johnson, Athanasios Konstantinidis, \textbf{M. H. Langston}, Janice O. McMahon, Benoit J. Meister, Paul D. Mountcastle, Aale Naqvi, Benoit Pradelle, Tahina Ramananandro, Sanket Tavarageri, Richard A. Lethin}{US 11,500,557 B2 \textperiodcentered{} Granted November 15, 2022}{}
\pub{\href{https://patents.google.com/patent/US11068178B2/}{Systems and Methods for Minimizing Communications}}{Muthu M. Baskaran, Thomas Henretty, Ann Johnson, Athanasios Konstantinidis, \textbf{M. H. Langston}, Janice O. McMahon, Benoit J. Meister, Paul D. Mountcastle, Aale Naqvi, Benoit Pradelle, Tahina Ramananandro, Sanket Tavarageri, Richard A. Lethin}{US 11,068,178 B2 \textperiodcentered{} Granted July 20, 2021}{}
\pub{\href{https://patents.google.com/patent/US10564949B2/}{System and Method for Generation of Event Driven, Tuple-Space Based Programs}}{Muthu M. Baskaran, Thomas Henretty, \textbf{M. H. Langston}, Richard A. Lethin, Benoit J. Meister, Nicolas T. Vasilache, David E. Wohlford}{US 10,564,949 B2 \textperiodcentered{} Granted February 18, 2020}{}
\pub{\href{https://patents.google.com/patent/US10540107B2/}{Systems and Methods for Energy Proportional Scheduling}}{Muthu M. Baskaran, Thomas Henretty, Ann Johnson, Athanasios Konstantinidis, \textbf{M. H. Langston}, Janice O. McMahon, Benoit J. Meister, Paul D. Mountcastle, Aale Naqvi, Benoit Pradelle, Tahina Ramananandro, Sanket Tavarageri, Richard A. Lethin}{US 10,540,107 B2 \textperiodcentered{} Granted January 21, 2020}{}
\pub{\href{https://patents.google.com/patent/US10496304B2/}{Systems and Methods for Minimizing Communications}}{Muthu M. Baskaran, Thomas Henretty, Ann Johnson, Athanasios Konstantinidis, \textbf{M. H. Langston}, Janice O. McMahon, Benoit J. Meister, Paul D. Mountcastle, Aale Naqvi, Benoit Pradelle, Tahina Ramananandro, Sanket Tavarageri, Richard A. Lethin}{US 10,496,304 B2 \textperiodcentered{} Granted December 3, 2019}{}
\pub{\href{https://patentimages.storage.googleapis.com/5a/eb/44/ce2fa33fda01d1/US10466349.pdf}{Systems and Methods for Efficient Targeting}}{Muthu M. Baskaran, Thomas Henretty, Ann Johnson, Athanasios Konstantinidis, \textbf{M. H. Langston}, Janice O. McMahon, Benoit J. Meister, Paul D. Mountcastle, Aale Naqvi, Benoit Pradelle, Tahina Ramananandro, Sanket Tavarageri, Richard A. Lethin}{US 10,466,349 B2 \textperiodcentered{} Granted November 5, 2019}{}
\pub{\href{https://patentimages.storage.googleapis.com/24/01/e0/614b430190ff34/US10209971.pdf}{Systems and Methods for Approximation Based Optimization of Data Processors}}{Muthu M. Baskaran, Thomas Henretty, Ann Johnson, Athanasios Konstantinidis, \textbf{M. H. Langston}, Richard A. Lethin, Janice O. McMahon, Benoit J. Meister, Paul Mountcastle}{US 10,209,971 B2 \textperiodcentered{} Granted February 19, 2019}{}
\pub{\href{https://patents.google.com/patent/US10180828B2/}{Systems and methods for power optimization of processors}}{Muthu M. Baskaran, Thomas Henretty, Ann Johnson, Athanasios Konstantinidis, \textbf{M. H. Langston}, Richard A. Lethin, Janice O. McMahon, Benoit J. Meister, Paul Mountcastle, Benoit Pradelle}{US 10,180,828 B2 \textperiodcentered{} Granted January 15, 2019}{}
\pub{\href{https://patents.google.com/patent/US10145871B2/}{Systems and methods for joint angle-frequency determination}}{Muthu M. Baskaran, Thomas Henretty, Ann Johnson, \textbf{M. H. Langston}, Richard A. Lethin, Janice O. McMahon, Benoit J. Meister, Paul Mountcastle}{US 10,145,871 B2 \textperiodcentered{} Granted December 4, 2018}{}
\pub{\href{https://patents.google.com/patent/US10095434B2/}{Systems and Methods for Efficient Determination of Task Dependences After Loop Tiling}}{Muthu M. Baskaran, Thomas Henretty, Ann Johnson, Athanasios Konstantinidis, \textbf{M. H. Langston}, Janice O. McMahon, Benoit J. Meister, Paul D. Mountcastle, Aale Naqvi, Benoit Pradelle, Tahina Ramananandro, Sanket Tavarageri, Richard A. Lethin}{US 10,095,434 B2 \textperiodcentered{} Granted October 9, 2018}{}
\pub{\href{https://patents.google.com/patent/US9471377B2/}{Systems and methods for parallelizing and optimizing sparse tensor computations}}{Muthu M. Baskaran, Thomas Henretty, \textbf{M. H. Langston}, Richard A. Lethin, Benoit J. Meister, Nicolas T. Vasilache}{US 9,471,377 B2 \textperiodcentered{} Granted October 18, 2016}{}
\subhead{Published applications}
\pub{\href{https://patents.google.com/patent/US20260050766A1/}{Efficient Attention in Transformer Neural Networks Using State Space Models}}{Mukul Gagrani, Junyoung Park, Raghavv Goel, Dalton James Jones, Wonseok Jeon, Matthew James Morse, \textbf{Matthew Harper Langston}, Mingu Lee, Christopher Lott}{US 2026/0050766 A1 \textperiodcentered{} Published February 19, 2026}{}
\pub{\href{https://patents.google.com/patent/US20260044449A1/}{Efficient Machine Learning Caching via Attention Output-Based Token Eviction}}{Raghavv Goel, Mukul Gagrani, Junyoung Park, Dalton James Jones, Mingu Lee, Wonseok Jeon, Matthew James Morse, \textbf{Matthew Harper Langston}, Christopher Lott}{US 2026/0044449 A1 \textperiodcentered{} Published February 12, 2026}{}
\pub{\href{https://patents.google.com/patent/WO2026035335A1/}{Memory-Constrained Attention in Machine Learning}}{Dalton James Jones, Junyoung Park, Matthew James Morse, Raghavv Goel, Mukul Gagrani, Mingu Lee, \textbf{Matthew Harper Langston}, Pierre-David Letourneau, Christopher Lott}{WO 2026/035335 A1 \textperiodcentered{} Published February 12, 2026}{}
\pub{\href{https://patents.google.com/patent/US20250315651A1/}{Polynomial Based Transformer}}{Pierre-David Letourneau, Manish Kumar Singh, Hong Cai, \textbf{Matthew Harper Langston}, Dalton James Jones, Fatih Murat Porikli}{US 2025/0315651 A1 \textperiodcentered{} Published October 9, 2025}{}
\pub{\href{https://patents.google.com/patent/US20250086431A1/}{Accelerating Complex Workloads Using Sparse Fast Chebyshev Interpolation}}{Pierre-David Letourneau, Dalton James Jones, Matthew James Morse, \textbf{Matthew Harper Langston}}{US 2025/0086431 A1 \textperiodcentered{} Published March 13, 2025}{}
\pub{\href{https://patents.google.com/patent/US20240354363A1/}{Identifying Quadratic Programming Solutions}}{Pierre-David Letourneau, Rania Hassen, Gary McGrath, \textbf{Matthew Harper Langston}}{US 2024/0354363 A1 \textperiodcentered{} Published October 24, 2024}{}
\pub{\href{https://patents.google.com/patent/US20230244935A1/}{Training of Neural Network with Polynomial Solver}}{Pierre-David Letourneau, \textbf{Matthew Harper Langston}, Richard Lethin, Matthew James Morse}{US 2023/0244935 A1 \textperiodcentered{} Published August 3, 2023}{}
\pub{\href{https://patents.google.com/patent/US20220374757A1/}{Efficient Computation of Hamiltonians for Ising Processors}}{Pierre-David Letourneau, \textbf{Matthew Harper Langston}, Noah Isaac Amsel, Richard Lethin}{US 2022/0374757 A1 \textperiodcentered{} Published November 24, 2022}{}
\pub{\href{https://patents.google.com/patent/US20220374687A1/}{Fast-Multidimensional Global Polynomial Solver (FM-GPS)}}{Pierre-David Letourneau, \textbf{Matthew Harper Langston}, Richard Lethin, Matthew James Morse}{US 2022/0374687 A1 \textperiodcentered{} Published November 24, 2022}{}


\section{Selected Invited Talks and Posters}
\begin{itemize}
  \item NeurIPS 2026, OPT 2026: Optimization for Machine Learning Workshop (Sydney): posters, \textit{Graph Coarsening Enables Escaping Local Optima in Compute-Constrained Nonconvex MAX-CUT} and \textit{Multiplicative vs.\ Additive State Coupling in Learned Optimizer Schedules}.
  \item MIT CSAIL, Cambridge, MA, Nov 2026: \textit{\href{https://www.csail.mit.edu/event/qualcomm-csail-ai-research-intelligence-edge}{Qualcomm at CSAIL: From AI Research to Intelligence at the Edge}}.
  \item ICLR 2026, Workshop on Memory for LLM-Based Agentic Systems (MemAgents), 2026: poster, \textit{CAOTE: Optimizing KV Cache Memory Through Attention Output Error-based Token Eviction}.
  \item ICLR 2026: poster, \textit{QUOKA: Query-Oriented KV Selection for Efficient LLM Prefill}.
  \item NeurIPS 2025: talk, \textit{PADRe: A Unifying Polynomial Attention Drop-in Replacement for Efficient Vision Transformer}.
  \item DARPA ITA3 Program Review, Orlando, FL, 2018: \textit{PARallel Fast Algorithms for Imaging Through Everything (PARFAITE)}.
  \item New York Scientific Data Summit (NYSDS), New York, NY, 2016: \textit{A Tool for Exploring Parallel Software Transformations}.
  \item DOE Office of Nuclear Physics SBIR/STTR Exchange Meeting, Gaithersburg, MD, 2016: \textit{Polyhedral Mapping Assistant and Visualizer (PUMA-V)}.
  \item IEEE VISSOFT, Bremen, Germany, 2015: \textit{Polyhedral User Mapping and Assistant Visualizer Tool for the R-Stream Auto-Parallelizing Compiler}.
  \item DARPA PERFECT PI Annual Meeting, Arlington, VA: \textit{Multi-Threaded Low-Communication FFTs on Heterogeneous Architectures (CPU + GPU)}.
  \item IEEE HPEC, Waltham, MA, 2013: \textit{Re-Introduction of Communication-Avoiding FMM-Accelerated FFTs with GPU Acceleration}.
  \item NSF RTG Symposium, Courant Institute, New York, NY, Sep 2012: \textit{An Adaptive FMM-Based PDE Solver in Three Dimensions for Various Boundary Conditions and Geometries}.
  \item NASA Langley Research Center, Safety Critical Avionics Systems Branch, Hampton, VA, Mar 2012: \textit{An Adaptive FMM-Based Free-Space PDE Solver in Three Dimensions}.
  \item Center for Plasma Edge Simulation Annual Conference, Georgia Tech, Atlanta, GA, Dec 2008: \textit{Further Developments in FMM Solvers}.
  \item Frontiers in Applied and Computational Mathematics, NJIT, Newark, NJ, May 2008: \textit{Fast Elliptic PDE Solver for Non-Homogeneous Force Distributions in Complex Geometries}.
  \item Princeton Plasma Physics Laboratory, Computational Plasma Physics Seminar, Princeton, NJ, Mar 2008: \textit{A Kernel-Independent FMM-Based Elliptic PDE Solver in Complex Geometries with Non-homogeneous Force Distributions}.
  \item Center for Plasma Edge Simulation Annual Conference, ORNL, Oak Ridge, TN, Dec 2007: \textit{FMM-Based Solvers}.
\end{itemize}

\section{Teaching}
\begin{itemize}
  \item Adjunct Assistant Professor, Courant Institute, NYU: graduate courses in Mathematics and Computer Science, 2013 -- \todo{end}. \todo{list courses}
  \item Adjunct Instructor, Graduate Discrete Mathematics, Courant Institute, NYU, Summers 2004--2008.
  \item Teaching Assistant, Courant Institute, NYU: Graduate Computer Graphics (Ken Perlin), Fall 2004;
        Graduate Experiments in Motion Capture (Chris Bregler), Spring 2004; Undergraduate Numerical Computing
        (Olof Widlund), Fall 2002.
  \item Recitation Leader, Undergraduate Algorithms (Allan Gottlieb), Courant Institute, NYU, Fall 2003.
\end{itemize}

\section{Professional Service}
\begin{itemize}
  \item Referee, SIAM Journal on Scientific Computing (SISC). \todo{add recent reviewing, e.g.\ ICLR, ICML, NeurIPS}
  \item Reviewer, SC08 and SC09 (International Conference for High Performance Computing, Networking, Storage
        and Analysis); SIGGRAPH 2004--2007.
  \item Contributor, DOE ASCAC Subcommittee Report \textit{Top Ten Exascale Research Challenges}, 2014.
  \item Text reviewer: K.~Ferland, \textit{Discrete Mathematics: An Introduction to Proofs and Combinatorics},
        1st ed., Houghton Mifflin, 2009; G.~Chartrand, A.~D.~Polimeni and P.~Zhang, \textit{Mathematical Proofs:
        A Transition to Advanced Mathematics}, 2nd ed., Addison Wesley, 2007.
\end{itemize}

\end{document}
